\subsection{Feature Name}
\label{sec:impl_feature_name}

% --- MOTIVATION (1-3 sentences) ---
% Always present. Reference SRs, mention parent UR in parentheses.
% Connect to the dependency chain: what previously introduced 
% features does this build on?

This section describes the implementation of ..., which addresses
SR\_XX\_YY and SR\_XX\_ZZ (derived from UR\_XX).
It builds on the component abstraction introduced in 
Section~\ref{sec:impl_component}.

% --- THEORY / BACKGROUND (0.5-1 page, optional) ---
% Include only if the implementation requires concepts not covered 
% in Chapter 2. Examples: IJCSA iteration scheme, bisection-based 
% event localization, zero-crossing functions.
% Skip entirely for features where the implementation IS the 
% explanation (e.g., port system, history recording).
% Reference Chapter 2 for shared foundations rather than repeating.

% --- IMPLEMENTATION (1-2 pages) ---
% The core of the subsection. Use the most appropriate medium:
% - Class diagram: when the contribution is structural 
%   (e.g., CoSimComponent hierarchy, Algorithm strategy pattern)
% - Code listing: when a specific algorithm or mechanism matters
%   (e.g., IJCSA loop, event bisection, state transfer)
% - Prose with inline code: for simpler features where a diagram 
%   or listing would be overkill
% Keep listings short (15-25 lines). Simplify or excerpt rather 
% than showing full source code.

% --- VERIFICATION (0.5-1 page) ---
% Show ONE well-chosen minimal example that demonstrates 
% correctness of this specific feature in isolation.
% Present: setup, expected behavior, observed result, brief 
% discussion.
% Use a figure (plot) or table for quantitative results.
% Note known limitations if any.
% Do NOT show comprehensive test coverage — that's not the 
% purpose. The case study (Ch.6) validates at system level.
